\section{Introduction}
Cloud applications depend on relationships among compute, storage, identity, messaging, and delivery services. Provisioning an individual resource is therefore only one part of deployment: identifiers must flow between resources, prerequisites must become ready, and later operations must preserve the distinction between intended configuration and observed state. Infrastructure-as-code systems express these requirements through machine-readable deployment descriptions, while visual modeling tools expose their structure directly. Winery and OpenTOSCA demonstrate that browser-based topology design and topology-driven execution are established approaches, rather than recent consequences of generative AI~\cite{winery2013,opentosca2013}. More recently, language-model research has examined whether natural-language requirements can produce infrastructure artifacts that satisfy explicit intent and actually deploy~\cite{iaceval2024,iacgen2026}. Together, these developments motivate interfaces that connect assistance to an inspectable execution representation.

The practical difficulty is coordinating representations and decisions across that interface. A user may draw a valid-looking connection that the application cannot implement, provide a configuration reference without its dependency edge, or accept an AI response that contains plausible but unavailable identifiers. Errors can arise before any cloud request is made, yet ordinary JSON validation cannot express every graph invariant. Conversely, a graph can satisfy all supported structural rules and still lack an execution role, a deployable function package, or a policy required by the intended workload. An application must communicate which checks it performs and where user decisions remain necessary. Otherwise, successful parsing, visual neatness, and a positive deployment receipt can become misleading substitutes for evidence about infrastructure behavior.

Existing work provides strong foundations but addresses different units of analysis. Modeling languages support richer representations and formal properties; IaC benchmarks measure generated code against requirements or deployment outcomes; commercial platforms combine visual components, contracts, and automated workflows~\cite{doml2024,multi2025,massdriver}. CloudCanvas does not resolve the limitations of these approaches generally. The narrower gap investigated here is how a single application-level graph contract can connect manual editing, structured AI drafts, persistence, and a bounded SDK execution path, and what additional structural faults its layers detect. This formulation avoids comparing a prototype's feature inventory with another system's measured generation scores. It also makes the evaluation feasible without claiming that a local graph experiment measures provider permissions, operational safety, or user productivity.

We present CloudCanvas as a concrete AWS-focused implementation of this workflow. Users create resource nodes manually or obtain an AI blueprint, inspect its configuration, and explicitly apply it to a saved sketch. Publishing is a separate action. Express controllers then perform ownership and deployment checks, build an execution order, resolve source outputs, and invoke resource-specific AWS SDK handlers. PostgreSQL records sketches, conversations, configurations, resource identifiers, statuses, and deployment outcomes. The architecture deliberately constrains the supported representation: eleven resource types and four attachment pairings are implemented. This bounded scope enables explicit rules but limits the architectures that can be expressed. It does not provide unrestricted cloud automation or a new infrastructure language standard.

The paper addresses two research questions. \textbf{RQ1:} On a declared synthetic conformance corpus, how much does the complete graph contract improve rejection of structural faults over schema-only validation and a schema-plus-topology ablation? \textbf{RQ2:} What local processing cost accompanies these checks, and which deliberately constructed boundary conditions remain outside their acceptance predicate? The first question concerns conformity to implemented representation rules; the second examines cost and limits. Neither asks whether the language models generate better architectures than competing tools. We make the following concrete contributions:
\begin{itemize}
\item An implementation description connecting an editable AWS graph, reviewed AI proposals, persisted state, and dependency-aware SDK execution.
\item A formal account of the shared contract's structural acceptance boundary, including the separation between reference binding and runtime output resolution.
\item A reproducible 72-instance conformance experiment, three independent process runs, two reduced-validation comparisons, and six boundary controls, with raw results and vector figures.
\end{itemize}
The remaining sections review related work, formulate the problem, explain the implementation and experiment, analyze the results, and discuss the limits of the supported claims.
